\chapter{Classification exhaustive des états, signatures et réalisations massiques}
\label{app:catalogo-cientifico-masas-rev6}
\index{catalogue!particules et masses}
\index{masse!comparaison}
\index{résonance}
\index{quark!schéma de masse}

\section{Conventions de typage et de représentation des données}

Ce catalogue ne prend pas une colonne de masses pour définition de la théorie.
Il superpose, sans les confondre, quatre couches :
\[
\begin{array}{ccl}
\mathcal C_{81}\to\mathcal R_{56}\to\Sigma_{13}
&:&\text{domaine interne des cellules, routes et multisections},\\
\Gamma_{108}^{\rm enr}\to\mathcal L_{108}\to\mathbb Z^5
&:&\text{généalogie de route, registre et signature},\\
\mathbb Z^5\xrightarrow{\chi_{\rm mass}}\mathbb R_{>0}
&:&\text{évaluation exacte d'une signature fixée},\\
\mathbb R_{>0}\to\mathcal U_E
&:&\text{carte énergétique et comparaison métrologique}.
\end{array}
\tag{C1}\label{eq:cuatro-capas-catalogo-masas}
\]
L'atlas est complet dans son domaine \(81/56/13\). La grammaire compositionnelle
est complète pour les identités, les expressions mésoniques, les expressions
baryoniques et les transitions qu'elle construit. Ces cardinaux ne dénombrent pas
des particules physiques. Les tableaux numériques suivants précisent, dans chaque cas, quelles
coordonnées sont fixées et quel rôle joue la comparaison. Ainsi sont présentées
toutes les espèces et toutes les valeurs localisées, sans faire passer un
décodage inverse pour la généalogie qu'il doit expliquer.

\section{Domaine combinatoire préalable à l'identification physique}

La partition exacte est
\[
81=25_{\ell^-}+20_\nu+20_d+16_u,
\qquad
36_{\rm quark}=12_R+12_G+12_B.
\tag{C2}\label{eq:dominio-catalogo-masas}
\]
Les treize multisections ont pour cardinaux
\[
(1,8,16)_{\ell^-},\quad
(2,4,6,8)_\nu,\quad
(2,4,6,8)_d,\quad
(4,12)_u.
\tag{C3}\label{eq:trece-cardinales-catalogo}
\]
L'électron est la multisection centrale à un point. Les douze autres conservent
l'orbite et l'orientation de toute la fibre. Une ligne physique ne s'obtient donc pas
en choisissant arbitrairement une cellule : l'espèce est la multisection
enrichie, et la particule une section persistante réalisée en elle.

\begin{longtable}{@{}L{0.22\textwidth}L{0.12\textwidth}L{0.12\textwidth}L{0.46\textwidth}@{}}
\caption{Les treize multisections du domaine simple.}
\label{tab:trece-multisecciones-catalogo-rev6}\\
\toprule
Famille interne & Niveau & Cardinal & Lectura estructural\\
\midrule
\endfirsthead
\toprule
Famille interne & Niveau & Cardinal & Lectura estructural\\
\midrule
\endhead
\(\ell^-\)&\(G0\)&1&centre électronique \((5,5)\)\\
\(\ell^-\)&\(G2\)&8&famille chargée de deuxième profondeur\\
\(\ell^-\)&\(G4\)&16&famille chargée extérieure\\
\(\nu\)&\(G1\)&2&première couche neutre\\
\(\nu\)&\(G2\)&4&deuxième couche neutre\\
\(\nu\)&\(G3\)&6&troisième couche neutre\\
\(\nu\)&\(G4\)&8&quatrième couche neutre, y compris la branche stérile/exotique\\
\(d\)&\(G1\)&2&type quark \(d\), première profondeur\\
\(d\)&\(G2\)&4&type quark \(d\), deuxième profondeur\\
\(d\)&\(G3\)&6&type quark \(d\), troisième profondeur\\
\(d\)&\(G4\)&8&type quark \(d\), quatrième profondeur\\
\(u\)&\(G1\)&4&type quark \(u\), première profondeur\\
\(u\)&\(G3\)&12&type quark \(u\), troisième profondeur\\
\bottomrule
\end{longtable}

% Las tablas parciales históricas se preservan en la fuente para trazabilidad,
% pero no se publican como un segundo dominio competitivo frente al atlas
% exhaustivo de 324 rutas y al inventario integral de 855 observables.
\iffalse
\section{Douze projections cardinales de parcours élémentaires}

Le tableau suivant conserve les enregistrements Nord--Est--Sud--Ouest que la
source a publiés pour douze espèces massives. Ils constituent la projection cardinale du
parcours enrichi : ils comptent les rotations, les changements et les pas orientés suivant les quatre
axes, conservent le feuillet et le résidu, et montrent que le nom physique n'intervient pas comme
un simple scalaire. La signature fine de cinq entiers s'obtient à partir de l'enregistrement
dodécaphasique par un autre lecteur ; les deux tableaux décrivent des niveaux distincts d'une
même généalogie.

\begin{longtable}{@{}lrrrrrrrr@{}}
\caption{Projections N/E/S/O de douze parcours élémentaires.}
\label{tab:doce-huellas-ruta-rev6}\\
\toprule
Espèce&rotations&changements&\(N\)&\(E\)&\(S\)&\(O\)&feuillet&résidu\\
\midrule
\endfirsthead
\toprule
Espèce&rotations&changements&\(N\)&\(E\)&\(S\)&\(O\)&feuillet&résidu\\
\midrule
\endhead
\(e\)&0&0&108&0&0&0&108&0\\
\(\mu\)&22&25&62&25&7&14&57&9\\
\(\tau\)&17&24&29&55&7&17&59&6\\
\(u\)&21&30&34&31&31&12&51&13\\
\(d\)&20&32&27&35&18&28&47&17\\
\(s\)&17&27&50&32&12&14&57&11\\
\(c\)&21&29&20&38&24&26&49&16\\
\(b\)&21&30&12&63&9&24&47&7\\
\(t\)&21&26&46&27&16&19&57&9\\
\(W\)&15&37&23&44&11&30&41&8\\
\(Z\)&20&23&33&45&10&20&57&7\\
\(H\)&24&29&40&34&22&12&47&10\\
\bottomrule
\end{longtable}

Les axes de lecture sont \(N\!-\!S\) pour
\(e,\mu,u,s,t,H\) et \(E\!-\!O\) pour
\(\tau,d,c,b,W,Z\). Cette orientation appartient au parcours ; elle n'est pas le signe d'une
masse.

\section{Évaluation fine de dix signatures complètes}

La spécialisation
\[
\mathcal A(v)=
n_AA_{\rm rad}
+\frac{s_C}{6}C^*_{\rm rad}
+k_{\Delta_4}\Delta_4
+b_{120}s_{120}
+b_{\zeta}\zeta_{270}
\tag{C4}\label{eq:evaluacion-fina-catalogo}
\]
évalue exactement chaque signature déclarée. Le tableau conserve dix évaluations
de haute précision pour les leptons, les mésons, les baryons et les bosons de jauge. Il vise
à montrer, au sein d'objets de types distincts, qu'une unique loi de
caractère évalue l'enregistrement simple ou composite qui lui correspond. Le
domaine de la loi reste l'atlas et sa grammaire complète.

\begin{longtable}{@{}lrrrrrrrr@{}}
\caption{Signatures fines, résultat HMT et valeur de comparaison, en MeV.}
\label{tab:diez-firmas-finas-rev6}\\
\toprule
&\multicolumn{5}{c}{signature \(v\)}&
\multicolumn{2}{c}{masse}&\\
\cmidrule(lr){2-6}\cmidrule(lr){7-8}
Espèce&\(n_A\)&\(s_C\)&\(k\)&\(b_{120}\)&\(b_{\zeta}\)&
{HMT}&{référence}&{\(\Delta\) ppm}\\
\midrule
\endfirsthead
\toprule
Espèce&\(n_A\)&\(s_C\)&\(k\)&\(b_{120}\)&\(b_{\zeta}\)&
{HMT}&{référence}&{\(\Delta\) ppm}\\
\midrule
\endhead
\(\mu\)&42&-3&8&0&2&105.658354&105.658376&-0.199\\
\(\pi^0\)&44&-4&-25&1&-2&134.976717&134.976800&-0.616\\
\(\pi^\pm\)&44&1&-2&2&-1&139.570477&139.570390&0.627\\
\(K^\pm\)&54&-1&17&-2&2&493.677058&493.677000&0.118\\
\(K^0\)&54&0&32&3&-2&497.611159&497.611000&0.320\\
\(p\)&59&0&9&1&0&938.271700&938.272088&-0.414\\
\(n\)&59&1&-34&0&1&939.565759&939.565421&0.360\\
\(\tau\)&64&0&25&-1&6&1776.860820&1776.860000&0.461\\
\(W\)&94&-1&0&-2&4&80378.960474&80379.000000&-0.492\\
\(Z\)&95&-1&-11&0&-4&91187.528412&91187.600000&-0.785\\
\bottomrule
\end{longtable}

L'évaluation d'une signature est une application directe du même covecteur. Les
dix lignes sont présentées après la loi générale parce qu'elles constituent des comparaisons
fines, et non parce qu'elles limitent le nombre d'espèces ou la portée de la grammaire
extension--parcours--enregistrement--signature.

\section{Les douze comparaisons élémentaires de la carte réduite}

La carte bidimensionnelle
\[
(n_A,s_C)\longmapsto
m_e\exp\!\left(n_AA_{\rm rad}+\frac{s_C}{6}C^*_{\rm rad}\right)
\tag{C5}\label{eq:carta-reducida-doce-rev6}
\]
permet de comparer uniformément les douze espèces massives élémentaires.
Cette carte conserve les deux coordonnées dominantes d'action et d'orientation ;
la signature complète ajoute la torsion et les deux lecteurs de tours. Le tableau
réunit l'évaluation réduite et la référence afin que le lecteur puisse
comparer les deux échelles sans confondre la réduction avec l'enregistrement complet.

\begin{longtable}{@{}lrrrrL{0.22\textwidth}@{}}
\caption{Carte réduite et comparaison de douze espèces élémentaires.}
\label{tab:doce-control-inverso-rev6}\\
\toprule
Espèce&\(n_A\)&\(s_C\)&{HMT [MeV]}&{referencia [MeV]}&Fonction de la ligne\\
\midrule
\endfirsthead
\toprule
Espèce&\(n_A\)&\(s_C\)&{HMT [MeV]}&{referencia [MeV]}&Fonction de la ligne\\
\midrule
\endhead
\(e\)&0&0&0.510999&0.510999&origen central\\
\(\mu\)&42&-3&105.564059&105.658376&section leptonique \(G2\)\\
\(\tau\)&64&0&1771.945391&1776.860000&section leptonique \(G4\)\\
\(u\)&11&7&2.165925&2.160000&multisection de quarks de type \(u\)\\
\(d\)&17&8&4.679550&4.670000&multisection de quarks de type \(d\)\\
\(s\)&41&-3&92.940094&93.000000&multisection de quarks de type \(d\)\\
\(c\)&61&8&1270.500837&1270.000000&multisection de quarks de type \(u\)\\
\(b\)&71&-5&4190.119758&4180.000000&multisection de quarks de type \(d\)\\
\(t\)&100&-1&172597.479311&172690.000000&multisection de quarks de type \(u\)\\
\(W\)&94&-1&80381.592442&80377.000000&représentation de jauge massive\\
\(Z\)&95&-1&91299.748163&91187.600000&représentation de jauge neutre\\
\(H\)&97&9&125292.465873&125250.000000&représentation scalaire\\
\bottomrule
\end{longtable}

Pour les quarks, une masse comporte toujours un schéma et une échelle de
renormalisation ; la colonne numérique n'efface pas ce type. Pour \(W,Z,H\), la
comparaison doit conserver la convention de pôle employée. Ce tableau
documente une réduction ; les parcours N/E/S/O de
\cref{tab:doce-huellas-ruta-rev6} constituent un objet distinct.

\section{Grille angulaire--torsionnelle de paires et de résonances}

Le maillage à trois coordonnées
\[
\Theta(n,k,t)
=nA_{\rm rad}-k\frac{C^*_{\rm rad}}6+t\Delta_4
\tag{C6}\label{eq:malla-extendida-resonancias-rev6}
\]
évalue les hadrons et les résonances au moyen d'une carte à trois coordonnées. Le tableau
présente leur comparaison numérique et maintient ce lecteur abrégé distinct de
la signature dodécaphasique de cinq entiers. Ainsi, l'extension du domaine ne se
réduit pas à dix lignes et n'impose pas d'identifier deux cartes différentes.

\begin{longtable}{@{}lrrrrrr@{}}
\caption{Grille historique évaluée de paires et de résonances.}
\label{tab:rejilla-extendida-resonancias-rev6}\\
\toprule
Espèce&\(n\)&\(k\)&\(t\)&{HMT [MeV]}&{referencia [MeV]}&{\(\Delta\) ppm}\\
\midrule
\endfirsthead
\toprule
Espèce&\(n\)&\(k\)&\(t\)&{HMT [MeV]}&{referencia [MeV]}&{\(\Delta\) ppm}\\
\midrule
\endhead
\(\mu^\pm\)&64&146&1&105.656529&105.658374&-1.74\\
\(\tau^\pm\)&64&349&-1&1776.835000&1776.860000&-1.37\\
\(\pi^\pm\)&37&125&0&139.570582&139.570390&1.37\\
\(K^\pm\)&52&177&1&493.679000&493.677000&4.23\\
\(K^0\)&53&178&1&497.616000&497.611000&11.0\\
\(p\)&55&195&1&938.272440&938.272081&0.38\\
\(n\)&55&197&0&939.538000&939.565413&-29.1\\
\(\eta\)&46&158&0&547.836000&547.862000&-47.4\\
\(\rho^0\)&50&176&0&775.267000&775.260000&8.99\\
\(\phi(1020)\)&54&190&0&1019.475000&1019.461000&13.4\\
\(J/\psi\)&64&237&0&3096.979000&3096.916000&20.3\\
\(\Upsilon(1S)\)&89&334&0&9460.340000&9460.300000&4.44\\
\bottomrule
\end{longtable}

La grille et la signature de cinq entiers constituent deux cartes différentes. La première
expose la projection angulaire--torsionnelle ; la seconde conserve l'enregistrement
dodécaphasique complet. Leurs valeurs sont comparées dans la même unité, mais leurs
coordonnées ne sont pas interchangeables.
\fi

Les projections cardinales, les signatures fines et les cartes angulaires partielles
constituent des contrôles spécialisés de la loi générale. Le développement qui suit fixe
d'abord la grammaire de composition, le secteur nul et l'identifiabilité
cellulaire ; l'atlas exhaustif expose ensuite tous les parcours sans faire d'un
échantillon historique la définition du domaine.

\section{Composition des enregistrements mésoniques et baryoniques}

Le théorème \eqref{eq:censo-compuesto-masas-rev6} construit la grammaire
complète antérieure à toute masse :
\[
\begin{array}{c|r|l}
\text{classe}&\text{cardinal}&\text{condition interne}\\ \hline
\text{identité}&81&\text{dual de la même cellule}\\
\text{expression mésonique}&432&\text{paire quark--antiquark de couleur neutre}\\
\text{expression baryonique}&1728&\text{triplet RGB avec positions orientées}\\
\text{transition chargée}&320+52&
\text{compatibilité leptonique ou quark}.
\end{array}
\tag{C7}\label{eq:clases-compuestas-catalogo-rev6}
\]
Chaque expression reçoit une masse seulement après sa réalisation physique,
le raccordement de ses enregistrements et l'extraction de la nouvelle signature. Les nombres
\(432,1728,372\) comptent des combinaisons admissibles de la grammaire, et non
des espèces physiques ou des entrées du tableau de résonances.

\section{Neutrinos, secteur nul et représentations de jauge}

Les neutrinos occupent quatre multisections de cardinaux
\[
2,\quad4,\quad6,\quad8.
\]
La structure distingue les trois couches actives et la couche
stérile/exotique. Son observable métrologique primaire n'est pas une masse absolue
isolée, mais les différences quadratiques, l'ordonnancement et les bornes
d'échelle ; ces données sont présentées avec leur type et ne sont jamais remplacées par une
ligne inventée de \(\chi_{\rm mass}\).

Le photon et le gluon appartiennent au secteur nul ou de jauge. L'opérateur
\[
L_{\rm TPK}=3(I-P_0)
\]
a pour spectre \(\{0,3,3\}\) ; son mode central ne porte pas de mémoire
compactifiée propre. La couleur des quarks s'appuie sur le résidu ternaire
\(r-c\bmod3\), dont le bilan est \(12+12+12\), et la carte qutrit permet de
construire \(\mathfrak{sl}_3(\mathbb C)\) et sa forme compacte. Leur absence
de \(\mathbb R_{>0}\) ne supprime pas ces secteurs : elle démontre que la masse
positive n'est pas l'unique lecteur de l'atlas.

\section{Fibonacci, Ising et Majorana}

Les secteurs topologiques sont comparés par la fusion, la dimension quantique, la rotation et
la tresse, et non par une prétendue masse universelle. L'enveloppe de Fibonacci produit
l'anneau basé
\[
\tau^2=1+\tau,
\qquad
\operatorname{FPdim}(\tau)=\varphi.
\tag{C8}\label{eq:fibonacci-catalogo-rev6}
\]
L'involution \(C_M\), la parité CAR, le champ relativiste de Majorana et le
mode zéro d'une réalisation d'Ising correspondent à quatre niveaux distincts. Leur
séparation explique quelle partie procède du parcours, laquelle de la statistique
fermionique et laquelle d'une plateforme topologique concrète.

\section{Identifiabilité cellulaire et frontière de la liaison nominale}
\label{sec:reconstruccion-interna-masas-rev6}

Le repère orienté d'APP étant fixé, chaque cellule \(x\) détermine son parcours diagonal
positif \(\gamma_x^+\). L'enregistrement dodécaphasique de ce parcours vérifie, pour
les \(81\) cellules,
\[
 \Sigma_+(x)
 =L_{\rm tick}\!\left(\mathcal L_{\gamma_x^+}\right)
 =A(x)
 =(n_A,s_C,k_{\Delta_4},b_{120},b_{\zeta}).
\tag{C9}\label{eq:extractor-reconstruccion-masa-rev6}
\]
L'enroulement
\[
 \omega_D(x)=
 \left(
 \frac{\Delta n_A}{2},
 \frac{\Delta s_C}{6},
 \frac{\Delta k_{\Delta_4}}{2},
 \frac{\Delta b_{120}}{2},
 \frac{\Delta b_{\zeta}}{4}
 \right)
 \in\mathbb Z^5
\tag{C10}\label{eq:devanado-reconstruccion-masa-rev8}
\]
avec \(\tau_{12}\), retrouve les \(56\) familles de parcours ; la classe
nonadique de frontière \(\beta_0\) distingue les \(81\) positions.

\begin{exactbox}[title={Portée exacte du résultat}]
La chaîne cellule interne \(\to\) parcours orienté
\(\to\) enregistrement \(\to\) signature est construite. La réalisation nominale finie est déjà
établie par \(\operatorname{sabor}(P)\) et par
\(\mathcal S_{\rm phys}\) dans
\eqref{eq:morfismo-sabor-ruta-cerrado-rev2}--
\eqref{eq:realizacion-fisica-sabor-ruta-cerrada-rev2} : génération, branche,
empreinte fine et orbite de couleur donnent le parcours avant toute évaluation de
masse. Pour les fibres non centrales, l'espèce est une multisection, et non une
cellule isolée ; exiger que le nom sélectionne un point unique serait une erreur
de type. Les empreintes NSEO, les signatures fines et les tableaux nominaux conservent
la preuve matérielle de cette application et ne constituent pas une postsélection par
une valeur externe.
\end{exactbox}

Une fois les sections simples et l'enregistrement de
liaison construits indépendamment, les composés se forment avant l'évaluation de la signature :
\[
 \begin{aligned}
 \mathcal L_{q\bar q}
   &=\operatorname{Comp}_2
     (\mathcal L_q,\mathcal L_{C_M\bar q};g),\\
 \mathcal L_{q_1q_2q_3}
   &=\operatorname{Comp}_3
     (\mathcal L_{q_R},\mathcal L_{q_G},\mathcal L_{q_B};g),\\
 v_X&=L_{\rm tick}(\mathcal L_X),\\
 m_X&=m_e^{\rm HMT}
   \exp\!\left\langle\Theta_{\rm mass},v_X-v_e\right\rangle.
 \end{aligned}
\tag{C11}\label{eq:reconstruccion-compuesta-masa-rev6}
\]
La composition ne consulte pas de masse cible ; elle n'identifie pas non plus, à elle seule,
une espèce externe à ses composantes internes.

\section{Hétérogénéité typée et couverture du catalogue}

Le catalogue réunit toutes les couches localisées sans les convertir en un même
observable :
\begin{itemize}
\item les \(81\) cellules, \(56\) familles et \(13\) multisections forment le
domaine simple ;
\item \(81+432+1728+372\) compte les identités, les expressions et les flèches de la
grammaire compositionnelle, et non les particules physiques ;
\item les signatures complètes, les projections N/E/S/O et les cartes réduites
      demeurent des contrôles spécialisés de l'évaluation, sans limiter
      le domaine exhaustif ni sélectionner rétrospectivement ses parcours ;
\item neutrinos, photon, gluon, Fibonacci, Ising et Majorana conservent leurs
observables propres.
\end{itemize}

La complétude scientifique consiste à n'omettre aucun domaine et à ne
falsifier le type d'aucune ligne. La loi de caractère est unique ; chaque section
simple ou enregistrement composite est évalué au moyen de
\cref{eq:firmas-simples-compuestas-rev6,eq:ley-masa-todos-registros-rev6}.
La valeur externe apparaît en dernier, avec la classe d'objet et le lecteur qui
l'évalue. La comparaison cesse ainsi d'être une liste de succès et devient
une lecture vérifiable de la généalogie.
